% titel: Blatt 6, A7: VSEPR – Struktur, Winkel, Hybridisierung
% skript: Mortimer Kap. 9 (chemie: VSEPR-Modell, Molekülgeometrie)

Das VSEPR-Modell sagt die Geometrie eines Moleküls aus der Zahl der
Elektronenpaare um das Zentralatom (Bindungs- und freie Paare) voraus, die sich
maximal abstossen. Aus der Paarzahl folgen Anordnung, Bindungswinkel und
Hybridisierung.

\begin{werkzeugkasten}[Elektronenpaar-Anordnung (EP-Paare $\to$ Form)]
Die Elektronenpaar-Anordnung bestimmt die Richtung der σ-Bindungen:
\begin{center}
\begin{tabular}{@{}c c c@{}}
  EP-Paare & Anordnung & Hybridisierung \\
  \hline
  2 & linear ($180^\circ$)          & $sp$  \\
  3 & trigonal-planar ($120^\circ$) & $sp^2$ \\
  4 & tetraedrisch ($109.5^\circ$)  & $sp^3$ \\
\end{tabular}
\end{center}
Freie Paare ''recken'' die tatsächliche Geometrie: aus 4 EP-Paaren mit einem
freien Paar wird trigonal-pyramidale; mit zwei einseitige (''V''/ecke) Form.
\textit{Skript:} Mortimer Kap. 9 (VSEPR, Hybridisierung).
\end{werkzeugkasten}

\begin{achtung}[Darauf achten]
\begin{itemize}
  \item Freie Elektronenpaare stossen stärker ab als Bindungen $\to$
        Bindungswinkel werden aus den idealen Werten \textit{komprimiert}
        (z.\ B. $107^\circ$ statt $109.5^\circ$ bei $\mathrm{NH_3}$, $104.5^\circ$
        bei $\mathrm{H_2O}$).
  \item Hybridisierung richtet sich nach der \textit{Elektronenpaar-Anordnung},
        nicht nach der Zahl blosser Atomnachbarn.
  \item Doppelbindungen zählen als \textit{ein} Elektronenpaar-Gebiet
        (bei $\mathrm{CO_2}$: 2 Gebiete, linear, $sp$).
\end{itemize}
\end{achtung}
